\documentclass[11pt]{article}
\usepackage[margin=1in]{geometry}
\usepackage{amsmath,amssymb,amsthm,mathtools}
\usepackage{microtype}
\usepackage[colorlinks=true,linkcolor=blue,urlcolor=blue]{hyperref}
\usepackage{booktabs,array}
\hypersetup{pdftitle={An all-threshold proof of Lemma H.1},
  pdfsubject={Actual facets, continuity, and Fatou's lemma}}
\newcommand{\R}{\mathbb R}
\newcommand{\E}{\mathbb E}
\newcommand{\PP}{\mathbb P}
\newcommand{\ind}{\mathbf 1}
\newcommand{\dd}{\,\mathrm d}
\newcommand{\ip}[2]{\langle #1,#2\rangle}
\newcommand{\norm}[1]{\lVert #1\rVert}
\newcommand{\Jf}{J_{\mathrm f}}
\newcommand{\Je}{J_{\mathrm e}}
\newcommand{\HH}{\mathcal H}
\newtheorem{lemma}{Lemma}
\newtheorem*{Hone}{Lemma H.1}
\title{An all-threshold proof of Lemma H.1\\[4pt]
  \large Actual facets, continuity, and Fatou's lemma}
\author{}
\date{}
\begin{document}
\maketitle
\vspace{-1.5em}

This note explains the proof used in the checked formalization of the
two localized Gaussian surface estimates. The facet estimate is proved
directly. For the ridge estimate, we first obtain a uniform bound at
almost every threshold, then pass to an arbitrary permitted threshold
by continuity and Fatou's lemma. The point of the last step is that
several facets may meet at a codimension-three face. We do not need a
pointwise Gauss--Green formula at such a nonsimple configuration.

\section{The statement and its geometric quantities}

Let $p\ge1$, set $L:=\log(ep)$, and let $u_1,\ldots,u_m\in\R^r$ be
the prescribed unit face normals, where $m\le2p$. Write
\[
 P(a):=\{x\in\R^r:u_j\cdot x\le a_j\text{ for every }j\},
 \qquad a\in\R^m.
\]
The notation $P(a)^s$ means $P(a+s\mathbf 1)$; in particular, these
are layers obtained by moving all supporting thresholds by the same
amount. Suppose that a random vector $V$ satisfies
\begin{equation}\label{eq:concentration}
 \PP\{V\in P(b)^s\setminus P(b)^{-s}\}
 \le C_0\kappa(s+\delta),
 \qquad b\in\R^m,\quad s>0,
\end{equation}
where $C_0>0$, $\kappa\ge0$, and $\delta\ge0$. No density of the law
of $V$ is assumed. Below $G\sim N(0,I_r)$ is independent of $V$.

For a proper full-dimensional polyhedron $P$, an \emph{actual facet}
is a face of dimension $r-1$. If its supporting equation is
$u_j\cdot g=b_j$, put
\begin{align}
 \Jf(P)&:=\sum_{\text{actual facets }F_j}|b_j|
       \int_{F_j}(\norm g^2-r)\varphi_r(g)\dd\HH^{r-1}(g),
       \label{eq:Jf}\\
 \Je(P)&:=\sum_{\text{actual ridges }E}\vartheta_E
       \int_E(\norm g^2-r)\varphi_r(g)\dd\HH^{r-2}(g).
       \label{eq:Je}
\end{align}
Here $\varphi_r(g):=(2\pi)^{-r/2}e^{-\norm g^2/2}$, a ridge is a
face of codimension two, and $\vartheta_E$ is the angle between the
outward unit normals of its two incident facets. Surface measure is
Euclidean surface measure, equivalently Lebesgue measure in an
orthonormal affine chart. We set both quantities to zero for the empty
set and the whole space. There are no ridge terms in dimension one.

\begin{Hone}
Suppose that \eqref{eq:concentration} holds and that $Q:=P(a)$ is empty
or full-dimensional. Then, for $0<\sigma\le1$,
\begin{equation}\label{eq:H1}
 \E\bigl|\Jf((Q-V)/\sigma)\bigr|
 +\E\bigl|\Je((Q-V)/\sigma)\bigr|
 \le C\kappa\bigl(\sigma L^{3/2}+\delta L^2\bigr),
\end{equation}
and both random quantities are integrable. The constant $C$ depends
only on $C_0$. No lower bound on the angles is required.
\end{Hone}

We prove the statement for each fixed target $a$. An important order
of choices is that a facet description of $Q$ will be fixed using only
$Q$ and the original inequalities. It will not depend on $V$, its
law, or $\sigma$.

\section{Fixing the actual facets}

If $Q$ is empty or the whole space, translation and positive scaling
preserve that property, so both terms on the left of \eqref{eq:H1}
vanish. Assume henceforth that $Q$ is proper and has nonempty interior.
Remove redundant inequalities, leaving
an index set $I\subseteq\{1,\ldots,m\}$ such that
\begin{equation}\label{eq:subsystem}
 Q=\{x:u_i\cdot x\le a_i,\ i\in I\},
\end{equation}
and no inequality in this description can be removed. After
reindexing, we use $i=1,\ldots,q$ for these retained inequalities.
They define the actual facets of $Q$.

There is a point $x_0$ with $u_i\cdot x_0<a_i$ for every retained $i$.
For each $i$, irredundancy also gives a point $y_i$ satisfying all the
other inequalities and violating the $i$th one. Along the segment
$(1-t)x_0+ty_i$, the $i$th inequality becomes equality at some
$0<t<1$. Every other inequality is still strict there, because it is
strict at $x_0$ and satisfied at $y_i$. This gives a point $x_i$ with
\begin{equation}\label{eq:facet-witness}
 u_i\cdot x_i=a_i,
 \qquad u_k\cdot x_i<a_k\quad(k\ne i).
\end{equation}
Thus each retained hyperplane supports a genuine facet. These strict
witnesses also show that the same retained inequalities are actual
facets at every sufficiently small perturbation of $a$: move $x_i$
slightly in direction $u_i$ to match the perturbed supporting level;
all the finitely many strict inequalities remain strict.

The concentration estimate survives this selection with the same
coefficient $C_0\kappa$ and the same $\delta$. To see this, assign
arbitrary thresholds to the retained inequalities and assign the
common threshold $M$ to the unused ones. As $M\to\infty$, the
indicators of the resulting symmetric face layers converge pointwise
to the indicator for the retained subsystem. Indeed, for each fixed
point, every unused inequality and its inward shift are eventually
satisfied. Bounded convergence applied to
\eqref{eq:concentration} proves the assertion. We may therefore work
with the fixed irredundant description \eqref{eq:subsystem}.

We record the geometric fact that will later make continuity possible.
For distinct retained indices $i,j,k$ with $u_i,u_j$ independent,
\begin{equation}\label{eq:no-affine-triple}
 u_k=\alpha u_i+\beta u_j
 \quad\Longrightarrow\quad
 a_k\ne\alpha a_i+\beta a_j.
\end{equation}
Suppose, to the contrary, that the supporting levels satisfy this
equality. If both coefficients are nonpositive, the positive slacks
at $x_0$ give
\[
 a_k-u_k\cdot x_0
 =\alpha(a_i-u_i\cdot x_0)+\beta(a_j-u_j\cdot x_0)\le0,
\]
a contradiction. If at least one coefficient is positive, rearranging
the relation, when necessary, expresses one of these three
inequalities as a nonnegative linear combination of the other two.
That inequality is redundant, again a contradiction. In particular,
the affine intersection of two independent retained supporting
hyperplanes is never contained in a third retained supporting
hyperplane. This still allows four or more facets to meet at a
codimension-three face.

\section{The uniform estimate at generic thresholds}

We explain the estimate that will be passed to the limit. For the
fixed retained normal family, let
\[
 R(g):=1+\max_S\norm{\Pi_Sg},
\]
where $S$ ranges over spans of at most three independent retained
normals, and $\Pi_S$ is orthogonal projection. If this family is empty,
the maximum is zero. There are at most $q+q^2+q^3$ such spaces.
Each projected norm is $1$-Lipschitz, so $R$ is $1$-Lipschitz and
$R\ge1$. A standard Gaussian projected onto a space of dimension at
most three satisfies
\[
 \E\exp\{\norm{\Pi_SG}^2/4\}\le2^{3/2}.
\]
The union bound therefore gives
\[
 \PP\{R(G)>1+t\}
 \le\min\{1,2^{3/2}(q+q^2+q^3)e^{-t^2/4}\}.
\]
Integrating this tail, splitting at a constant multiple of
$\sqrt L$, proves
\begin{equation}\label{eq:moments}
 \E R(G)^k\le C_kL^{k/2},\qquad k=1,2,3,4.
\end{equation}
The constants are independent of the ambient dimension.

For every linearly dependent subcollection of distinct normals, fix
one nonzero relation $\sum_i c_i u_i=0$. Call a threshold vector
\emph{generic} if it avoids the corresponding equations
$\sum_i c_i b_i=0$. There are only finitely many subcollections, and
each excluded equation is a proper hyperplane in threshold space.
Thus the generic set has full Lebesgue measure. At a generic
threshold, every simultaneously active set of normals is independent.
In particular a nonempty polyhedron is full-dimensional: at any one
of its points, choose a direction whose inner product with each
active normal is negative, and move a sufficiently small distance in
that direction to obtain a strict interior point. An actual
codimension-three face has exactly three incident facets. Only this
three-facet case is needed in the ridge calculation below.

For a deterministic polyhedron $P$ at such a threshold, the geometric
estimate is
\begin{equation}\label{eq:localization}
 |\Jf(P)|+|\Je(P)|
 \le C\int R(g)^4
   \ind\{g\in P^{1/R(g)}\setminus P\}\varphi_r(g)\dd g.
\end{equation}
We give the reasons for this estimate, since they explain both the
powers of $L$ and the absence of an angle-gap assumption.

On a facet, write $g=b_iu_i+y$ with $y\perp u_i$. The tangential
divergence of $y\varphi_r(g)$ is
$(r-1-\norm y^2)\varphi_r(g)$. Integrating by parts on every facet and
pairing the contributions of neighboring facets yields
\begin{equation}\label{eq:facet-green}
 \Jf(P)=\sum_i|b_i|(b_i^2-1)\int_{F_i}\varphi_r
 -\sum_{E=F_i\cap F_j}
   \bigl(|b_i|d_{j\mid i}+|b_j|d_{i\mid j}\bigr)
       \int_E\varphi_r.
\end{equation}
Here $n_{j\mid i}$ is the outward conormal to $E$ within $F_i$ and
$d_{j\mid i}=g\cdot n_{j\mid i}$ on $E$. If $z_E$ is the projection
of a point of $E$ onto its two-dimensional normal space, elementary
two-normal geometry gives
\begin{equation}\label{eq:pair-cancellation}
 \bigl||b_i|d_{j\mid i}+|b_j|d_{i\mid j}\bigr|
 \le2\vartheta_E\norm{z_E}^2.
\end{equation}
For completeness, write the coefficient as
$|b_i|(d_{j\mid i}+d_{i\mid j})+(|b_j|-|b_i|)d_{i\mid j}$.
The bounds $|b_i|,|d_{i\mid j}|\le\norm{z_E}$,
$|d_{j\mid i}+d_{i\mid j}|\le\vartheta_E\norm{z_E}$, and
$\bigl||b_i|-|b_j|\bigr|\le\vartheta_E\norm{z_E}$ prove
\eqref{eq:pair-cancellation}. Small angles occur as multiplicative
weights; no inverse angle is introduced.

For the ridge quantity, integrate the tangential divergence once more,
now within each ridge. At a generic threshold this gives
\begin{equation}\label{eq:generic-edge-green}
 \Je(P)=\sum_E\vartheta_E(\norm{z_E}^2-2)\rho_E
 -\sum_H\ip{z_H}{\sum_{E\supset H}\vartheta_E n_{E,H}}\rho_H,
\end{equation}
where $H$ runs over actual codimension-three faces,
$\rho_F:=\int_F\varphi_r$, and $n_{E,H}$ is the outward conormal to
$H$ within $E$. The second sum is absent if $r\le2$.
The projection $z_H$ onto the normal space of $H$ is constant on $H$.
Its normal cone $C_H$ cuts out a spherical triangle $\Omega_H$.
In this three-dimensional normal space, the facet of $C_H$
corresponding to $E$ is a planar wedge of angle $\vartheta_E$ and has
outward normal $-n_{E,H}$. Gaussian integration by parts gives
\[
 \int_{C_H}x e^{-\norm x^2/2}\dd x
 =\sum_{E\supset H}\vartheta_E n_{E,H},
\]
because the Gaussian integral over a planar wedge is its angle:
its radial factor is $\int_0^\infty t e^{-t^2/2}\dd t=1$.
Polar coordinates express the same vector integral as
$2\int_{\Omega_H}\omega\dd\HH^2(\omega)$, since the radial factor is
$\int_0^\infty t^3e^{-t^2/2}\dd t=2$. Consequently
\[
 \sum_{E\supset H}\vartheta_E n_{E,H}
 =2\int_{\Omega_H}\omega\dd\HH^2(\omega),
 \qquad
 \left\|\sum_{E\supset H}\vartheta_E n_{E,H}\right\|
 \le2|\Omega_H|.
\]
This use of \eqref{eq:generic-edge-green} is restricted to generic
thresholds. The passage to a nonsimple target will use continuity of
$\Je$, rather than the same boundary formula at that target.

We next turn the surface expressions into a volume integral. For
$x$ in the relative interior of a face $F$ of codimension
$s\in\{1,2,3\}$, take
\[
 g=x+y,\qquad y\in N_F,\qquad
          0<\norm y<\frac1{4R(x)},
\]
where $N_F$ is the outward normal cone. These regions are disjoint
for distinct faces. Indeed, for $w\in P$ one has
$y\cdot(w-x)\le0$, and hence
\[
 \norm{g-w}^2\ge\norm y^2+\norm{x-w}^2.
\]
Thus $x$ is the unique closest point of $P$ to $g$. Distinct relative
face interiors cannot supply the same closest point.

The tangent and normal coordinates are orthogonal. Since the normal
space is spanned by at most three of the prescribed normals,
$|x\cdot y|\le R(x)\norm y\le1/4$. Consequently
\[
 \frac{\varphi_r(x+y)}{\varphi_r(x)}
 =e^{-x\cdot y-\norm y^2/2}
\]
is bounded above and below by numerical constants. The same is true
of $R(x+y)/R(x)$, by the Lipschitz property of $R$. The volumes of
the normal pieces are respectively
\[
 \frac1{4R(x)},\qquad
 \frac{\vartheta_E}{32R(x)^2},\qquad
 \frac{|\Omega_H|}{192R(x)^3}.
\]
Integrating $R(g)^4\varphi_r(g)$ over these regions therefore controls
the facet integrands $R(x)^3\varphi_r(x)$, the ridge integrands
$\vartheta_E R(x)^2\varphi_r(x)$, and the codimension-three
integrands $|\Omega_H|R(x)\varphi_r(x)$. These are exactly the
terms in \eqref{eq:facet-green}--\eqref{eq:generic-edge-green}.
Moreover, $g\notin P$, and
$u_i\cdot g\le b_i+\norm y$ for every $i$. The choice of radius gives
$\norm y R(g)<1$, so the regions lie in
$P^{1/R(g)}\setminus P$. Summing over the disjoint regions proves
\eqref{eq:localization}. Unbounded faces cause no difficulty: insert
smooth cutoffs at radius $M$ in the divergence calculations, then let
$M\to\infty$; Gaussian polynomial integrability makes the cutoff
terms tend to zero.

The facet part of this argument only uses the two-facet cancellation
\eqref{eq:pair-cancellation}. It therefore applies directly to the
fixed actual-facet presentation of $Q$, including when its
codimension-three faces are nonsimple. In particular it already gives
the required bound for $\Jf$ after averaging as follows.

For thresholds $b$ in the retained subsystem and $v\in\R^r$, put
\[
 P_{b,v}:=\{g:u_i\cdot g\le(b_i-u_i\cdot v)/\sigma,\ i\le q\}.
\]
The events and surface quantities here are measurable; the fixed-chart
formula \eqref{eq:fixed-chart} below expresses them as parameter
integrals of measurable Gaussian functions.
For each fixed $v$, the map
$b\mapsto(b_i-u_i\cdot v)_{i\le q}/\sigma$ is an invertible affine
map of threshold space and preserves null sets. Fubini's theorem
therefore shows that, for Lebesgue-almost every $b$, the generic
ridge calculation applies to $P_{b,V}$ almost surely, even if the law
of $V$ is singular.

To average \eqref{eq:localization}, condition on $G=g$. Its layer event
implies
\[
 V+\sigma g\in P(b)^{\sigma/R(g)}\setminus P(b).
\]
As an event in $V$, this is an outer face layer with thresholds
$b_i-\sigma u_i\cdot g$ and width $\sigma/R(g)$. It is contained in
the symmetric layer appearing in \eqref{eq:concentration}. Thus its
conditional probability is at most
$C_0\kappa(\sigma/R(g)+\delta)$. Tonelli's theorem and
\eqref{eq:moments} now give
\begin{equation}\label{eq:averaged-generic}
 \E|\Je(P_{b,V})|
 \le C\kappa\{\sigma\E R(G)^3+\delta\E R(G)^4\}
 \le C\kappa(\sigma L^{3/2}+\delta L^2)
 \quad\text{for almost every }b.
\end{equation}
The same averaging gives the facet bound at the fixed target $a$.
The right side is uniform in $b$. Its finiteness establishes
integrability at the good thresholds; integrability was not assumed
when applying Tonelli.

\section{Continuity at the fixed target}

We now prove the precise continuity needed to remove the exceptional
thresholds from \eqref{eq:averaged-generic}. Keep the selected actual
facets fixed. For every fixed $v\in\R^r$, we claim that
\begin{equation}\label{eq:continuity}
 b\longmapsto\Je(P_{b,v})
 \quad\text{is continuous at }b=a.
\end{equation}
The claim concerns the value of the surface integral. It does not
require differentiability of a threshold probability at $a$.

Consider two nonparallel retained normals $u_i,u_j$. Their normal
space $S_{ij}$ and its orthogonal complement $T_{ij}$ are fixed.
The simultaneous supporting equations have a unique normal-space
solution $z_{ij}(b,v)\in S_{ij}$; it is linear in the two supporting
levels, and hence continuous in $b$. Every point on the corresponding
affine plane is uniquely $z_{ij}(b,v)+y$, with $y\in T_{ij}$.
For $k\ne i,j$, let
\[
 w_k:=\Pi_{T_{ij}}u_k,\qquad
 d_k(b,v):=\frac{b_k-u_k\cdot v}{\sigma}
                         -u_k\cdot z_{ij}(b,v).
\]
The remaining inequalities on this plane are $w_k\cdot y\le d_k(b,v)$.

Using the fixed orthonormal tangent chart, the contribution of this
pair to $\Je$ is
\begin{multline}\label{eq:fixed-chart}
 \vartheta_{ij}\,(2\pi)^{-1}e^{-\norm{z_{ij}(b,v)}^2/2}
 \int_{T_{ij}}
   \bigl(\norm{z_{ij}(b,v)}^2+\norm y^2-r\bigr)\\[-3pt]
 \times\prod_{k\ne i,j}\ind\{w_k\cdot y\le d_k(b,v)\}
          \varphi_{T_{ij}}(y)\dd y.
\end{multline}
Here $\varphi_{T_{ij}}$ is the standard Gaussian density on the
$(r-2)$-dimensional space $T_{ij}$. A pair that does not bound an
actual ridge contributes zero: its section is empty or has smaller
dimension and therefore has zero tangent Gaussian measure. Summing
\eqref{eq:fixed-chart} over unordered nonparallel pairs gives the
actual ridge quantity. Parallel opposite facets do not intersect in a
full-dimensional polyhedron and supply no ridge term.

Let $b_n\to a$. If $w_k\ne0$, the equality
$w_k\cdot y=d_k(a,v)$ defines a Gaussian-null hyperplane in
$T_{ij}$. Off this hyperplane the corresponding indicators converge.
If $w_k=0$, then $u_k$ lies in the span of $u_i,u_j$. Relation
\eqref{eq:no-affine-triple}, preserved by translation and positive
scaling, says that $d_k(a,v)\ne0$. In this case the inequality is
constant in $y$ and its truth value is unchanged for all sufficiently
large $n$. There are only finitely many $k$, so the product of the
indicators converges for almost every tangent Gaussian point.

For fixed $v$, the vectors $z_{ij}(b_n,v)$ remain bounded. The absolute
integrand in \eqref{eq:fixed-chart} is therefore bounded by
$C_v(1+\norm y^2)\varphi_{T_{ij}}(y)$, an integrable function. Dominated
convergence proves continuity of the pair contribution. A finite sum
proves \eqref{eq:continuity}. The facet witnesses ensure that nearby
thresholds retain the chosen actual facets, so the pair sums used here
agree with the geometric quantities in a neighborhood of the target.

This argument permits many facets at a codimension-three face. For a
fixed pair chart, the additional supporting hyperplanes are genuine
lower-dimensional boundaries, so their intersections remain Gaussian
null. Most importantly, continuity holds for \emph{every fixed $v$}.
The null sets used above belong to the tangent Gaussian integration;
they do not require $V$ itself to avoid any set. This is why an atomic
or otherwise singular law of $V$ is allowed. The bound $C_v$ need not
be integrable in $v$: we will use Fatou, not dominated convergence,
for the outer expectation.

\section{Fatou's lemma and completion of the proof}

Fix the law of $V$ and $\sigma>0$, and write
\[
 B:=C\kappa(\sigma L^{3/2}+\delta L^2).
\]
Let $\mathcal G\subseteq\R^q$ be the full-measure set of thresholds
at which \eqref{eq:averaged-generic} and integrability hold. A
full-measure subset of Euclidean space is dense: every nonempty open
ball has positive Lebesgue measure and hence meets $\mathcal G$.
Choose $b_n\in\mathcal G$ with $b_n\to a$, staying in the actual-facet
neighborhood just described. The selected facet family stays fixed
throughout this choice. The sequence is allowed to depend on the
law of $V$ and on $\sigma$; the facet selection and $B$ are uniform in
the required data.

Set $Y_n:=|\Je(P_{b_n,V})|$ and $Y:=|\Je(P_{a,V})|$. These random
variables are measurable, since their fixed-chart formulas are
parameter integrals of measurable functions. By
\eqref{eq:continuity}, $Y_n\to Y$ pointwise in $V$. Apply Fatou's
lemma to these nonnegative functions:
\begin{equation}\label{eq:fatou}
 \E Y
 =\E\bigl[\liminf_{n\to\infty}Y_n\bigr]
 \le\liminf_{n\to\infty}\E Y_n
 \le B.
\end{equation}
At this stage the first expectation is an extended nonnegative
integral. Its finite upper bound proves $Y\in L^1$, and therefore
proves integrability of the signed ridge quantity itself. There is
no circular assumption of integrability at the target.

Because $P_{a,V}=(Q-V)/\sigma$, \eqref{eq:fatou} is the required
ridge estimate at the original, possibly nonsimple $Q$. Add the
direct facet estimate to obtain \eqref{eq:H1}, enlarging the constant
only by an absolute factor. The coefficient $C_0$ in the base
concentration estimate is absorbed into $C$. Every preceding step
uses only $\sigma>0$, so this proof in fact gives the stated estimate
for all positive $\sigma$; the paper needs only $\sigma\le1$.

\section*{Connection with the checked proof}

The note describes the continuity/Fatou route already checked in the
\href{https://github.com/VC2015/rectangle-clt-singular}{rectangle-CLT repository}
at proof commit
\begin{center}
\texttt{443fbc667ed3c3ee47369e028c2c4280a5eed604}.
\end{center}
The relevant Lean files are in the directory
\begin{center}
\small\path{lean-rectangle-clt-arxiv-v1/RectangleCLTArxivV1/}
\end{center}
\begin{center}
\small
\begin{tabular}{@{}p{0.38\textwidth}p{0.57\textwidth}@{}}
\toprule
Mathematical step & Lean source file\\
\midrule
Fixed actual facets & \texttt{SupplementalHActualFacetSelection.lean}\\
Inherited concentration & \texttt{SupplementalHSubsystemConcentration.lean}\\
Direct facet expectation & \texttt{SupplementalHActualFacetAveraged.lean}\\
Almost-everywhere ridge bound & \texttt{SupplementalHRidgeMain.lean}\\
Continuity of ridge values & \texttt{SupplementalHRidgeValueContinuity.lean}\\
Nonnegative Fatou argument & \texttt{SupplementalHSurfaceFatou.lean}\\
Fixed-threshold ridge bound & \texttt{SupplementalHActualRidgeEstimates.lean}\\
Final empty/full-dimensional theorem & \texttt{SupplementalHSurfaceFinal.lean}\\
\bottomrule
\end{tabular}
\end{center}
The endpoint is
\texttt{exists\_fixed\_surface\_bound\_of\_empty\_or\_fullDimensional}.
The fixed numerical constant is chosen before the model data. The
formal proof uses Euclidean orthonormal-chart surface measures.
The generic ridge calculation and the limiting argument are separate
checked components; this note does not claim a completed formalization
of the printed nonsimple ridge Gauss--Green identity.

\end{document}
